\documentclass[11pt,a4paper]{article}

\usepackage[margin=24mm]{geometry}
\usepackage[T1]{fontenc}
\usepackage{amsmath}
\usepackage{booktabs}
\usepackage{listings}
\usepackage[hidelinks]{hyperref}
\lstset{
    language=Python,
    basicstyle=\ttfamily\footnotesize,
    breaklines=true,
    columns=fullflexible,
    keepspaces=true,
    showstringspaces=false,
    frame=single
}
\setlength{\parskip}{4pt}
\setlength{\parindent}{0pt}

\title{Labwork 5 Report: Gaussian Blur and Shared Memory}
\author{Dao Quy Tung \\ Student ID: 2540109}
\date{3 October 2026}

\begin{document}
\maketitle

\section{Objective}
For this lab, I changed the grayscale kernel from Labwork 4 to apply a
$7\times7$ Gaussian blur. I wrote two versions: one reads the filter from
global memory, and the other copies the filter into shared memory.
I used \texttt{time.time()} to compare their running times.

\section{Implementation}
\subsection{Filter and pixel calculation}
I used the same $1200\times600$ RGB image,
\texttt{2021-Road-Race-Grand-Championship.jpg}. The filter contains the
following integer weights:
\[
K=
\begin{bmatrix}
0 & 0 & 1 & 2 & 1 & 0 & 0\\
0 & 3 & 13 & 22 & 13 & 3 & 0\\
1 & 13 & 59 & 97 & 59 & 13 & 1\\
2 & 22 & 97 & 159 & 97 & 22 & 2\\
1 & 13 & 59 & 97 & 59 & 13 & 1\\
0 & 3 & 13 & 22 & 13 & 3 & 0\\
0 & 0 & 1 & 2 & 1 & 0 & 0
\end{bmatrix}.
\]

Each thread calculates one output pixel. For each colour channel, it
multiplies the values in the surrounding $7\times7$ area by the filter
weights and adds them together. The weights sum to 1003, so the final
value is divided by 1003 using integer division:
\[
O_c(y,x)=
\left\lfloor
\frac{\displaystyle\sum_{j=0}^{6}\sum_{i=0}^{6}
K_{j,i}\,I_c(y+j-3,x+i-3)}{1003}
\right\rfloor.
\]
Here, $c$ is the red, green, or blue channel. If a neighbouring position
falls outside the image, $I_c$ uses the nearest edge pixel. In the code,
this is done by keeping the row and column coordinates within the valid
range. The result is stored as \texttt{uint8} in a separate output array.

Both kernels use the 2D indexing from Labwork 4, with
$16\times16$ blocks and a $75\times38$ grid. Threads outside the image
do not calculate an output pixel.

\subsection{Shared-memory version}
In \texttt{gaussian\_blur}, each thread reads the filter weights directly
from global memory. In \texttt{gaussian\_blur\_shared}, the threads in each
block first copy the 49 weights into a shared array:
\begin{lstlisting}
sharedFilter = cuda.shared.array((7, 7), dtype=int32)
tid = cuda.threadIdx.x + cuda.threadIdx.y * cuda.blockDim.x
threadsPerBlock = cuda.blockDim.x * cuda.blockDim.y
for k in range(tid, 49, threadsPerBlock):
    sharedFilter[k // 7, k % 7] = weights[k // 7, k % 7]

cuda.syncthreads()
\end{lstlisting}

The synchronization makes sure the filter has been fully copied before
any thread uses it. All threads reach this point, including those outside
the image; the pixel boundary check comes afterwards. The calculation
then uses \texttt{sharedFilter[j, i]} for each weight. The input image
stays in global memory, and the blur calculation is the same in both
versions.

\section{Timing procedure}
I transferred the image and filter to the GPU and allocated both output
arrays before measuring time. I then ran each kernel once and called
\texttt{cuda.synchronize()} to finish the warm-up. The timed sections are:
\begin{lstlisting}
start = time.time()
gaussian_blur[gridSize, blockSize](devInput, devOutput, devFilter)
cuda.synchronize()
gpuTime = time.time() - start

start = time.time()
gaussian_blur_shared[gridSize, blockSize](devInput, devOutputShared, devFilter)
cuda.synchronize()
sharedTime = time.time() - start
\end{lstlisting}

These times cover the kernel launch and the wait for it to finish.
Memory allocation, transfers between the CPU and GPU, initial
compilation, and saving the images are outside the timed sections.
Copying the filter from global to shared memory happens inside the
shared-memory kernel, so that step is included in its time.
I calculated speedup as
\[
S=\frac{T_{\mathrm{without\ shared}}}{T_{\mathrm{with\ shared}}}.
\]

\section{Measured results}
I ran \texttt{Labwork5.py} on an NVIDIA GeForce RTX 2050 on
3 October 2026. The results from this run were:
\begin{center}
\begin{tabular}{lr}
\toprule
Measurement & Result \\
\midrule
GPU without shared memory & 0.000797 s \\
GPU with shared memory & 0.000735 s \\
Speedup & $1.08\times$ \\
\texttt{np.array\_equal(hostOutput, hostOutputShared)} & \texttt{True} \\
\bottomrule
\end{tabular}
\end{center}

The shared-memory version was slightly faster in this run.
The output comparison returned \texttt{True}, so both versions produced
exactly the same values for all 720,000 pixels and all three channels.
The results are saved as \texttt{blur\_gpu.png} and
\texttt{blur\_gpu\_shared.png}.

The speedup uses the full timing values before they are rounded for
display. I measured each version once with $16\times16$ blocks, so this
result only describes this run; the times may change on another run.

\end{document}
